\originalchapter{球形坍缩、暗晕与质量函数}
本章对应8.902第15讲末与第16讲。线性增长在$|\delta|\ll1$时可靠，暗晕却处在高度非线性状态。均匀球形过密区提供一个能完整求解的桥梁：由其真实非线性演化确定坍缩时刻，再把早期扰动按线性理论外推到这个时刻，建立用于统计的阈值。阈值与真实密度必须始终分开。

\originalsection{均匀过密球的参数解}
先取Einstein--de Sitter背景、无压强尘埃、球对称且壳层尚未交叉。半径$R$内总质量$M$恒定，边缘满足
\[
 \ddot R=-\frac{GM}{R^2},\qquad \frac12\dot R^2-\frac{GM}{R}=E.
\]
过密增长模最终转向，取$E<0$、$A=GM/(2|E|)$。令$R=A(1-\cos\vartheta)$，代入能量式选膨胀分支，可得$\dd t=B(1-\cos\vartheta)\dd\vartheta$，其中$A^3=GMB^2$。因此
\begin{equation}\label{eq:astro-top-hat-param}
 R=A(1-\cos\vartheta),\qquad t=B(\vartheta-\sin\vartheta).
\end{equation}
把大爆炸时刻设为背景与扰动共同的$t=0$排除这一模型的衰减模；若另设不同起始时间，不能继续使用相同阈值。转向发生在$\vartheta=\pi$，$R_{\rm ta}=2A$、$t_{\rm ta}=\pi B$；形式坍缩在$\vartheta=2\pi$、$t_c=2t_{\rm ta}$。

背景$\bar\rho=1/(6\pi Gt^2)$，球内$\rho=3M/(4\pi R^3)$。用参数解与$A^3=GMB^2$消去维度常数得
\begin{equation}\label{eq:astro-nonlinear-density}
 1+\delta_{\rm nl}=\frac92\frac{(\vartheta-\sin\vartheta)^2}{(1-\cos\vartheta)^3}.
\end{equation}
原笔记归一变量$x=a/a_{\rm ta}$、$y=R/R_{\rm ta}$给相同结果$1+\delta=\xi(x/y)^3$，其中$\xi=9\pi^2/16$；个别原式漏掉立方。牛顿运动式中的$\xi$也不能在尚未用边界条件求出前直接替换为3。

\originalsection{1.686阈值是怎样算出的}
早期$\vartheta\ll1$时，展开
\[
 1-\cos\vartheta=\frac{\vartheta^2}{2}-\frac{\vartheta^4}{24}+\cdots,\quad
 \vartheta-\sin\vartheta=\frac{\vartheta^3}{6}-\frac{\vartheta^5}{120}+\cdots.
\]
代入式\eqref{eq:astro-nonlinear-density}得$\delta_{\rm nl}=3\vartheta^2/20+O(\vartheta^4)$。线性增长模$\delta_L\propto t^{2/3}$，用这个早期系数定其归一化后，任意参数时刻的线性外推值为
\[
 \delta_L(\vartheta)=\frac3{20}[6(\vartheta-\sin\vartheta)]^{2/3}.
\]
所以转向时$\delta_{L,\rm ta}=\frac35(3\pi/4)^{2/3}\simeq1.062$，形式坍缩时
\begin{equation}\label{eq:astro-delta-collapse}
 \delta_c=\frac35\left(\frac{3\pi}{2}\right)^{2/3}\simeq1.68647.
\end{equation}
与此相对，转向时真实密度比$1+\delta_{\rm nl}=9\pi^2/16\simeq5.55$，形式坍缩时则发散。\textbf{$\delta_c$不是暗晕的真实密度对比}，它是把早期同一扰动沿线性增长延伸所得的统计标签。

\originalsection{为什么估计的维里半径是转向半径的一半}
均匀球的势能$U=-3GM^2/(5R)$。转向时宏观径向动能为零，故$E=U_{\rm ta}$。理想化最终稳定暗晕若没有净能量流出且满足$2T+U=0$，则$E=U_{\rm vir}/2$，给
\[
 R_{\rm vir}=\frac12R_{\rm ta}.
\]
将维里化时刻近似设为形式坍缩时刻$2t_{\rm ta}$，背景密度已降低4倍，球的密度则比转向时增大8倍。因此
\[
 \Delta_{\rm vir}\equiv\frac{\rho_{\rm vir}}{\bar\rho(t_c)}
 =32\frac{9\pi^2}{16}=18\pi^2\simeq177.65.
\]
严格无压强均匀球本身在$R=0$奇点终止，\emph{并未自行产生维里化}。真实体系的壳层交叉、相混合与非球对称扰动提供松弛；上面的半径与时刻是能量加维里条件的模型估计。有限压强、表面压强、潮汐、质量交换与$\Lambda$会改变关系，不能把它当成任意暗晕定理。

常用$M_{200c}$定义平均密度为$200\rho_{\rm crit}$，$M_{200m}$则为$200\bar\rho_m$。Einstein--de Sitter里二者参考密度相同，其他背景下不同。给出$M_{200}$时须说明后缀与红移，再用$M=(4\pi/3)R^3\Delta\rho_{\rm ref}$换算，不能把200倍平均密度和200倍临界密度任意互换。

\originalsection{Gaussian尾部如何估计已坍缩的质量比例}
选定共动球窗半径$R(M)=[3M/(4\pi\bar\rho_{m,0})]^{1/3}$，将\emph{线性}密度场平滑到$M$尺度，方差$\sigma^2(M,z)=D^2(z)\sigma^2(M,0)$由式\eqref{eq:astro-smoothed-variance}计算。若初始场为Gaussian，平滑密度的概率密度为
\[
 p(\delta_M)=\frac1{\sqrt{2\pi}\sigma}\exp\left(-\frac{\delta_M^2}{2\sigma^2}\right).
\]
单点超过球形阈值的概率是
\[
 F_{\rm naive}(>M,z)=\int_{\delta_c}^{\infty}p(\delta_M)\dd\delta_M
 =\frac12\operatorname{erfc}\left(\frac{\delta_c}{\sqrt2\sigma(M,z)}\right).
\]
按初始近均匀的Lagrange区域质量来解释，这估计的是\emph{质量比例}；不是高度非均匀、已坍缩Euler空间中的几何体积分数。一个在小尺度平滑后低于阈值的点，仍可能属于更大尺度已经超过阈值的区域，这称云中之云问题。

当$M\to0$且$\sigma\to\infty$时，上式趋于$1/2$。Press--Schechter处方将它乘2，取$F=\operatorname{erfc}(\nu/\sqrt2)$，$\nu=\delta_c/\sigma$。对于锐$k$窗产生的Markov随机游走和恒定吸收边界，反射原理可导出同一因子；实空间球窗有相关步长时不再自动成立。扩展处方的具体边界与首次通过计算在公开PS6题解中另作推导。这一模型不表示所有暴胀都精确产生Gaussian场，也不意味着所有暗晕遵守单一球形阈值。

\originalsection{质量比例转为数密度时的符号与单位}
$F(>M)$随$M$增大而减小，所以$\dd F/\dd M<0$。质量在$[M,M+\dd M]$中的比例是$-F_M\dd M$，乘今日平均物质密度并除以对象质量，得\emph{共动}质量函数
\begin{equation}\label{eq:astro-press-schechter}
 \frac{\dd n_c}{\dd M}=\sqrt{\frac2\pi}\frac{\bar\rho_{m,0}}{M^2}
 \frac{\delta_c}{\sigma(M,z)}\left|\frac{\dd\ln\sigma}{\dd\ln M}\right|
 \exp\left[-\frac{\delta_c^2}{2\sigma^2(M,z)}\right].
\end{equation}
若用$\bar\rho_m(z)$代替$\bar\rho_{m,0}$，得到的是物理体积数密度，相差$(1+z)^3$。原笔记同时使用$\rho_c(z)\Omega_m(z)$和“共动”名称，需作这一转换；其$F_M$还漏负号。窗的方差中必须是$|\widetilde W|^2$并带$1/(2\pi^2)$，不能用单次$\widetilde W$或$1/(2\pi)$。

当小质量处$\sigma\to\infty$且大质量处$\sigma\to0$，积分$\int_0^\infty M(\dd n_c/\dd M)\dd M=\bar\rho_{m,0}$。有微观截止或不满足这些边界时，不能继续说该处方必然把所有质量放入暗晕。$\nu\gg1$的稀有高质量暗晕受到指数抑制，对增长因子和振幅很敏感；小质量部分的形状则还依赖$\sigma(M)$的斜率。

球形模型是解析基准。椭球坍缩、合并、环境与实际暗晕识别方法改变精确质量函数，课程展示的Sheth--Tormen、Jenkins、Tinker等曲线是各自条件下的改进或模拟标定，并非从本章简单积分直接推得。比较模型必须统一质量定义、背景宇宙学、红移和体积约定，才能判断真正的差异。
